% SYNTHETIC FIXTURE for number_ledger.py tests. Numbers are deliberately inconsistent. Not a real paper.
\documentclass{article}
\title{Evolutionary Search for Efficient Attention}
\begin{document}
\maketitle
\begin{abstract}
We run an evolutionary architecture search for $\sim$1 month and find a model that improves accuracy from 73.1 to 78.0, a 16\% improvement over the baseline \citep{smith2024}.
\end{abstract}

\section{Introduction}
Our search takes $\sim$1 month on a shared cluster. Prior work in 2023 used far less compute (Section~\ref{sec:method}).

\section{Method}
\label{sec:method}
The search ran for $\sim$3 weeks in total. In every round we keep $K = 16$ elites and run 6 iterations of mutation. The total compute was 160 GPUs for 21 hours. Of the 207 runs, 12.4\% of 207 diverged. The fingerprint has 12 + 6 + 3 = 21 features.

\begin{figure}[t]
\caption{Selection keeps the 12 elites with the highest fitness after each round (Figure 2).}
\end{figure}

\begin{table}[t]
\caption{Compute budget.}
\begin{tabular}{lc}
\toprule
Stage & GPU-hours \\
\midrule
Search & 2,400 \\
Evaluation & 480 \\
Total & 2,880 \\
\bottomrule
\end{tabular}
\end{table}

\section{Results}
\begin{table}[t]
\caption{Accuracy per benchmark.}
\begin{tabular}{lcccc}
\toprule
Method & A & B & C & Avg \\
\midrule
Baseline & 70.0 & 72.0 & 74.0 & 72.0 \\
Ours & 76.0 & 78.0 & 80.0 & 81.0 \\
\bottomrule
\end{tabular}
\end{table}

Our fitness score is 0.68 vs 0.42 for the baseline, a 3$\times$ gain of 4.5 points. In the ablation, the 4 + 10 + 7 = 21 features are split by type.

\section*{References}
Smith et al. 2024. Something with 99 pages and 3 volumes.
\end{document}
